\chapter[DONNÉES APRS DANS LE CHAMP INFORMATION AX.25]{Données APRS dans le champ Information
AX.25}\label{donnuxe9es-aprs-dans-le-champ-information-ax.25}

\begin{aprssideblock}{Format générique\\des données}
\phantomsection\addcontentsline{toc}{section}{Format générique des données}\label{format-guxe9nuxe9rique-des-donnuxe9es}
En général, le champ Information AX.25 peut contenir tout ou partie des
informations suivantes :

\begin{itemize}
\tightlist
\item
  APRS Data Type Identifier
\item
  APRS Data
\item
  APRS Data Extension
\item
  Comment
\end{itemize}
\end{aprssideblock}

\begin{aprsbodyblock}
\setlength{\arrayrulewidth}{0.75pt}
\begin{tabular}{|c|c|c|c|}
\hline
\multicolumn{4}{|>{\columncolor{gray!20}}l|}{\bfseries\itshape Generic APRS Information Field}\\
\hline
\bfseries\itshape Data Type ID & \bfseries\itshape APRS Data & \bfseries\itshape APRS Data Extension & \bfseries\itshape Comment \\
\hline
\makebox[0pt][r]{Octets :\hspace{3.2em}}1 & n & 7 & n \\
\hline
\end{tabular}
\end{aprsbodyblock}

\begin{aprssideblock}{APRS Data Type\\Identifier}
\phantomsection\addcontentsline{toc}{section}{APRS Data Type Identifier}\label{aprs-data-type-identifier}
Chaque paquet APRS contient un APRS Data Type Identifier (DTI). Celui-ci
détermine le format du reste des données du champ Information, comme suit :
\end{aprssideblock}

\vspace{8pt}
\noindent\textbf{APRS Data Type Identifiers}\par\vspace{4pt}
{\footnotesize
\setlength{\tabcolsep}{4pt}
\setlength{\arrayrulewidth}{0.75pt}
\renewcommand{\arraystretch}{1.15}
\noindent
\begin{tabular}[t]{|>{\columncolor{gray!20}\centering\arraybackslash}p{26pt}|p{\dimexpr0.49\textwidth-46pt\relax}|}
\hline
\rowcolor{gray!20}
\bfseries\itshape Ident & \multicolumn{1}{>{\columncolor{gray!20}}c|}{\bfseries\itshape Data Type} \\
\hline
\texttt{0x1c} & Current Mic-E Data (Rev 0 beta) \\
\hline
\texttt{0x1d} & Old Mic-E Data (Rev 0 beta) \\
\hline
! & Position sans timestamp (sans messagerie), ou Ultimeter 2000 WX Station \\
\hline
" & \emph{[Unused]} \\
\hline
\# & Peet Bros U-II Weather Station \\
\hline
\$ & Raw GPS data ou Ultimeter 2000 \\
\hline
\% & Agrelo DFJr / MicroFinder \\
\hline
\& & \emph{[Reserved --- Map Feature]} \\
\hline
\textquotesingle{} & Old Mic-E Data (mais \emph{Current} data pour TM-D700) \\
\hline
( & \emph{[Unused]} \\
\hline
) & Item \\
\hline
* & Peet Bros U-II Weather Station \\
\hline
+ & \emph{[Reserved --- Shelter data with time]} \\
\hline
, & Invalid data ou test data \\
\hline
- & \emph{[Unused]} \\
\hline
. & \emph{[Reserved --- Space weather]} \\
\hline
/ & Position avec timestamp (sans messagerie) \\
\hline
\texttt{0–9} & \emph{[Do not use]} \\
\hline
: & Message \\
\hline
; & Object \\
\hline
\end{tabular}
\hfill
\begin{tabular}[t]{|>{\columncolor{gray!20}\centering\arraybackslash}p{26pt}|p{\dimexpr0.49\textwidth-46pt\relax}|}
\hline
\rowcolor{gray!20}
\bfseries\itshape Ident & \multicolumn{1}{>{\columncolor{gray!20}}c|}{\bfseries\itshape Data Type} \\
\hline
\textless{} & Station Capabilities \\
\hline
= & Position sans timestamp (avec messagerie APRS) \\
\hline
\textgreater{} & Status \\
\hline
? & Query \\
\hline
@ & Position avec timestamp (avec messagerie APRS) \\
\hline
\texttt{A–S} & \emph{[Do not use]} \\
\hline
T & Telemetry data \\
\hline
\texttt{U–Z} & \emph{[Do not use]} \\
\hline
{[} & Maidenhead grid locator beacon (obsolète) \\
\hline
\textbackslash{} & \emph{[Unused]} \\
\hline
{]} & \emph{[Unused]} \\
\hline
\textasciicircum{} & \emph{[Unused]} \\
\hline
\_ & Weather Report (sans position) \\
\hline
\textasciigrave{} & Current Mic-E Data (\emph{non utilisé} sur TM-D700) \\
\hline
\texttt{a–z} & \emph{[Do not use]} \\
\hline
\{ & User-Defined APRS packet format \\
\hline
\textbar{} & \emph{[Do not use --- TNC stream switch character]} \\
\hline
\} & Third-party traffic \\
\hline
\textasciitilde{} & \emph{[Do not use --- TNC stream switch character]} \\
\hline
\end{tabular}
}

\begin{aprsbodyblock}
Note : la radio Kenwood TM-D700 utilise le DTI \lit{\textquotesingle{}} pour
les données Mic-E actuelles (\emph{current}). Cette radio n'utilise pas le
DTI \lit{\textasciigrave{}}.
\end{aprsbodyblock}

\begin{aprssideblock}{APRS Data and\\Data Extension}
\phantomsection\addcontentsline{toc}{section}{APRS Data and Data Extension}\label{aprs-data-and-data-extension}
Il existe 10 types principaux d'APRS Data :

\begin{itemize}
\tightlist
\item
  Position
\item
  Direction Finding
\item
  Objects and Items
\item
  Weather
\item
  Telemetry
\item
  Messages, Bulletins and Announcements
\item
  Queries
\item
  Responses
\item
  Status
\item
  Other
\end{itemize}

Certaines de ces données peuvent également comporter une APRS Data
Extension, qui fournit des informations supplémentaires.

L'APRS Data et l'APRS Data Extension optionnelle suivent le Data Type
Identifier.

Le tableau de la page suivante présente une liste complète de tous les
différents types possibles d'APRS Data et d'APRS Data Extension.
\end{aprssideblock}

\begin{center}
{\footnotesize
\setlength{\tabcolsep}{4pt}
\setlength{\arrayrulewidth}{0.75pt}
\begin{tabular}{|>{\columncolor{gray!20}\raggedleft\arraybackslash\bfseries}m{0.15\linewidth}|>{\raggedright\arraybackslash}p{0.38\linewidth}|>{\raggedright\arraybackslash}p{0.38\linewidth}|}
\hline
\rowcolor{gray!20}
 & \multicolumn{1}{c|}{\bfseries\itshape Possible APRS Data} & \multicolumn{1}{c|}{\bfseries\itshape Possible APRS Data Extension} \\
\hline

Position &
Time (DHM or HMS)\newline
Lat/long coordinates\newline
Compressed lat/long/course/speed/radio range/altitude\newline
Symbol Table ID and Symbol Code\newline
Mic-E longitude, speed and course, telemetry or status\newline
Raw GPS NMEA sentence\newline
Raw weather station data &
Course and Speed\newline
Power, Effective Antenna Height/Gain/Directivity\newline
Pre-Calculated Radio Range\newline
Omni DF Signal Strength\newline
Storm Data (in Comment field) \\
\hline
Direction Finding &
Time (DHM or HMS)\newline
Lat/long coordinates\newline
Compressed lat/long/course/speed/radio range/altitude\newline
Symbol Table ID and Symbol Code &
Course and Speed\newline
Power, Effective Antenna Height/Gain/Directivity\newline
Pre-Calculated Radio Range\newline
Omni DF Signal Strength\newline
Bearing and Number/Range/Quality (in Comment field) \\
\hline
Objects and Items &
Object name\newline
Item name\newline
Time (DHM or HMS)\newline
Lat/long coordinates\newline
Compressed lat/long/course/speed/radio range/altitude\newline
Symbol Table ID and Symbol Code\newline
Raw weather station data &
Course and Speed\newline
Power, Effective Antenna Height/Gain/Directivity\newline
Pre-Calculated Radio Range\newline
Omni DF Signal Strength\newline
Area Object\newline
Storm Data (in Comment field) \\
\hline
Weather &
Time (MDHM)\newline
Lat/long coordinates\newline
Compressed lat/long/course/speed/radio range/altitude\newline
Symbol Table ID and Symbol Code\newline
Raw weather station data &
Wind Direction and Speed\newline
Storm Data (in Comment field) \\
\hline
Telemetry &
Telemetry (non Mic-E) & \\
\hline
Messages, Bulletins and Announcements &
Addressee\newline
Message Text\newline
Message Identifier\newline
Message Acknowledgement\newline
Bulletin ID, Announcement ID\newline
Group Bulletin ID & \\
\hline
Queries &
Query Type\newline
Query Target Footprint\newline
Addressee (Directed Query) & \\
\hline
Responses &
Position\newline
Object/Item\newline
Weather\newline
Status\newline
Message\newline
Digipeater Trace\newline
Stations Heard\newline
Heard Statistics\newline
Station Capabilities &
Course and Speed\newline
Power, Effective Antenna Height/Gain/Directivity\newline
Pre-Calculated Radio Range\newline
Omni DF Signal Strength\newline
Area Object\newline
Wind Direction and Speed \\
\hline
Status &
Time (DHM zulu)\newline
Status text\newline
Meteor Scatter Beam Heading/Power\newline
Maidenhead Locator (Grid Square)\newline
Altitude (Mic-E)\newline
E-mail message & \\
\hline
Other &
Third-Party forwarding\newline
Invalid Data/Test Data & \\
\hline
\end{tabular}
}
\end{center}

\begin{aprssideblock}{Comment Field}
\phantomsection\addcontentsline{toc}{section}{Comment Field}\label{comment-field}
En général, tout paquet APRS peut contenir un commentaire en texte libre
(tel qu'un message de beacon) dans le champ Information, immédiatement après
les APRS Data ou l'APRS Data Extension.

Il n'y a pas de séparateur entre les données APRS et le commentaire, sauf
mention contraire.

Le commentaire peut contenir tout caractère ASCII imprimable ou UTF-8 (à
l'exception de \lit{\textbar{}} et \lit{\textasciitilde{}}, réservés au TNC
channel switching).

Ne placez aucun retour chariot (0x0d) ni saut de ligne (0x0a) à la fin. Les
stations IGate les supprimeront, ce qui produira un contenu légèrement
différent.

La longueur maximale, très faible et apparemment arbitraire, du champ
Comment dépend du rapport concerné --- les détails figurent dans la
description de chaque rapport.

Dans certains cas particuliers, le champ Comment peut également contenir des
données APRS supplémentaires :

\begin{itemize}
\tightlist
\item
  \textbf{Altitude} dans le texte de commentaire (voir chapitre 6 : Formats
  de temps et de position), ou dans le texte de statut Mic-E (voir chapitre
  10 : Format Mic-E).
\item
  \textbf{Maidenhead Locator} (grid square), dans un champ de texte de
  statut Mic-E (voir chapitre 10 : Format Mic-E) ou dans un Status Report
  (voir chapitre 16 : Rapports de statut).
\item
  Les paramètres \textbf{Bearing and Number/Range/Quality}
  (\texttt{/BRG/NRQ}), dans les DF reports (voir chapitre 7 : Extensions de
  données APRS).
\item
  \textbf{Area Object Line Widths} (voir chapitre 11 : Rapports d'Objets et
  d'Items).
\item
  \textbf{Signpost Objects} (voir chapitre 11 : Rapports d'Objets et
  d'Items).
\item
  \textbf{Weather and Storm Data} (voir chapitre 12 : Rapports météo).
\item
  \textbf{Beam Heading and Power}, dans les Status Reports (voir chapitre
  16 : Rapports de statut).
\end{itemize}
\end{aprssideblock}

\begin{aprssideblock}{Base-91\\Notation}
\phantomsection\addcontentsline{toc}{section}{Base-91 Notation}\label{notation-base-91}
Deux formats de données APRS utilisent la notation base-91 : les
coordonnées lat/long au format compressé (voir chapitre 9 : Formats
compressés de rapport de position) et l'altitude au format Mic-E (voir
chapitre 10 : Format Mic-E).

Les données base-91 sont compressées en une courte chaîne de caractères.
Tous les caractères sont des caractères ASCII imprimables, avec des codes de
caractère compris entre 33 et 124 en décimal (c'est-à-dire de \lit{!} à
\lit{\textbar{}}).

Pour calculer la chaîne de caractères ASCII base-91 correspondant à une
valeur de donnée donnée, cette valeur est divisée par des puissances de 91
progressivement décroissantes jusqu'à ce que le reste soit inférieur à 91. À
chaque étape, 33 est ajouté au modulo de la division pour obtenir le code de
caractère ASCII correspondant.

Par exemple, pour une valeur de donnée de 12345678 :

\begin{verbatim}
12345678 / 91³ = modulo 16, reste 288542
288542   / 91² = modulo 34, reste 6988
6988     / 91¹ = modulo 76, reste 72
\end{verbatim}

Les quatre codes de caractère ASCII sont donc 49 (c'est-à-dire 16+33), 67
(c'est-à-dire 34+33), 109 (c'est-à-dire 76+33) et 105 (c'est-à-dire 72+33),
correspondant à la chaîne ASCII \texttt{1Cmi}.
\end{aprssideblock}

\begin{aprssideblock}{APRS Precision\\and Datum Option}
\phantomsection\addcontentsline{toc}{section}{APRS Precision and Datum Option}\label{aprs-precision-and-datum-option}
Cette extension de la spécification APRS permet de transmettre une plus
grande précision entre utilisateurs. Mais comme la précision sans
l'exactitude du datum n'a pas de sens, ce format inclut également
l'identification du datum utilisé pour obtenir cette précision.

\textbf{DATUM :} le datum par défaut pour le GPS et pour APRS est WGS84. La
spécification APRS indique que toutes les positions on-air sont présumées
être en WGS84 sauf indication contraire. Cette \guillemotleft{}option\guillemotright{}
documente donc deux méthodes pour \guillemotleft{}indiquer autrement\guillemotright{} :

\textbf{PAR PAYS :} ou par accord local continental. Deux accords de ce type
ont été déclarés : WGS84 en Amérique du Nord et OSGB36 au Royaume-Uni. Si
d'autres peuvent être convenus, ils seront listés.

\textbf{PAR DATUM BYTE :} le DATUM byte dans cette \guillemotleft{}PRECISION
and DATUM Option\guillemotright{}.

\textbf{FORMAT :} le format de cette option APRS Precision-and-Datum est la
présence d'un champ de 5 octets \texttt{!DAO!} apparaissant n'importe où dans
le champ de commentaire de position. Cette option est rétrocompatible car le
format de base DDMM.mmN/DDDMM.mmH est conservé et sera toujours décodé par
toutes les applications existantes. Le \texttt{!DAO!} fournit simplement la
précision supplémentaire jusqu'à environ un pied (30~cm) et identifie
également le datum :

\begin{verbatim}
!DAO! - longueur fixe, n'importe où dans le commentaire de position
 D    - identifiant du datum (base-91)
  A   - précision supplémentaire de LAtitude (base-91)
   O  - précision supplémentaire de LOngitude (base-91)
\end{verbatim}

\textbf{RECOMMANDATION :} il est recommandé de placer cette option à la FIN
de tout autre texte de commentaire de position. Ainsi, elle ne déplace aucun
texte de commentaire lisible par un humain que l'on souhaiterait afficher sur
les anciens systèmes. Cette option ajoutée peut dépasser la limite
\guillemotleft{}visible\guillemotright{} actuelle de 57 octets de la
spécification, puisque ces octets supplémentaires ne sont de toute façon pas
utilisés par les systèmes existants.

\textbf{PRÉCISION :} cette option offre trois degrés de précision.

\begin{enumerate}
\item \textbf{LISIBLE HUMAINEMENT, millièmes de minute.} Précis à environ
  6 pieds (2~m) près et, étant lisible humainement, permet même aux
  utilisateurs de D7, D700 ou de toute autre application existante de lire
  une position à 3 décimales de minutes. Ce format est identifié si le Datum
  byte est en majuscule.

\item \textbf{BASE-91.} Ajoute une précision supplémentaire au
  91/10\,000\textsuperscript{e} de minute près, soit environ 4 chiffres
  fractionnaires décimaux de minute, soit environ un pied (30~cm). Ce format
  est identifié si le datum byte est en minuscule.

\item \textbf{NULL.} Si les octets A et O sont des caractères ESPACE, ils ne
  sont présents que pour respecter le format \texttt{!DAO!} et n'impliquent
  AUCUNE précision ajoutée. Utilisé lorsqu'on souhaite envoyer une
  information de DATUM sans revendiquer de précision supplémentaire. Cette
  utilisation de caractères espace pour indiquer l'absence de chiffres de
  précision est cohérente avec le système d'ambiguïté APRS existant.
\end{enumerate}

\textbf{DATUMS :} plusieurs catégories de datums peuvent être utilisées.

\textbf{PRÉDÉFINIS :} les lettres A--Z et a--z indiquent l'un des 26 datums
courants. La casse de ces 26 lettres indique le type de précision utilisé.
Les majuscules indiquent des chiffres décimaux lisibles pour A et O, et les
minuscules indiquent un encodage base-91 pour A et O. Il existe 26 DATUMS
prédéfinis. Une table des 26 datums courants sera préparée. Voici quelques
exemples :

\begin{verbatim}
W = WGS84
N = NAD27
O = OSGB36
\end{verbatim}

\textbf{LOCAL CUSTOM :} pour des événements spéciaux et fermés, le DATUM
peut être l'une des dix options définies localement. Elles sont indiquées par
les chiffres 0 à 9. Une utilisation créative consiste, pour des événements
fermés utilisant des cartes ou des lieux fixes, à définir les cartes par
numéro~!

\textbf{ENCODAGE BASE-91 :} le base-91 est fréquemment utilisé en APRS pour
améliorer la résolution d'un octet sans utiliser de binaire et en
n'utilisant que le sous-ensemble de caractères ASCII imprimables. L'usage
exclusif de caractères ASCII imprimables présente de nombreux avantages qui
ont été abondamment discutés par ailleurs. Le base-91 signifie simplement la
valeur d'un octet ASCII après soustraction de 33 en décimal. Ainsi le
caractère \guillemotleft{}\texttt{!}\guillemotright{} moins 33 donne
\guillemotleft{}0\guillemotright{}, jusqu'au caractère
\guillemotleft{}\texttt{\}}\guillemotright{} moins 33 qui donne la valeur
\guillemotleft{}90\guillemotright{}.

\textbf{AUTRES :} il reste 28 autres possibilités d'information de DATUM
dans le chiffre D qui peuvent être utilisées si nécessaire.

\textbf{Exemples :}

\texttt{!W23!} signifie WGS84 et la majuscule indique qu'il s'agit du
format lisible humainement au 3\textsuperscript{e} chiffre. \guillemotleft{}2\guillemotright{}
est le troisième chiffre décimal des minutes de latitude et est lisible
humainement. \guillemotleft{}3\guillemotright{} est le chiffre ajouté des
minutes de longitude et est également lisible humainement.

\texttt{!wAb!} signifie WGS84 mais la minuscule indique qu'il s'agit de la
précision ajoutée au pied près. \guillemotleft{}A\guillemotright{} est le
code base-91 pour deux chiffres supplémentaires (65 moins 33 donne
\guillemotleft{}32\guillemotright{}) et \guillemotleft{}b\guillemotright{}
donne deux chiffres supplémentaires de longitude (98 -- 33, soit
\guillemotleft{}...65\guillemotright{}).

\texttt{!w:\textbackslash!} serait également WGS84 mais avec
\guillemotleft{}:\guillemotright{} et \guillemotleft{}\textbackslash\guillemotright{}
décodant deux chiffres supplémentaires de \guillemotleft{}27\guillemotright{}
pour la latitude et \guillemotleft{}59\guillemotright{} pour la longitude,
au pied près environ.

\textbf{CONVERSION BASE-91 :} dans le premier exemple \texttt{!W23!}
ci-dessus, les chiffres réels sont simplement ajoutés aux LAT/LONG
existantes, par exemple pour ajouter de la précision à DDMM.mmN/DDDMM.mmH
et devenir équivalent à DDMM.mm2N et DDDMM.mm3H.

Mais dans les deux exemples suivants \texttt{!wAb!} et \texttt{!w:\textbackslash!},
les deux chiffres ajoutés ne peuvent pas être simplement ajoutés à la chaîne
de position ASCII, puisqu'ils ne peuvent aller que de 00 à 90. Il faut donc
les mettre à l'échelle pour qu'ils aillent de 00 à 99. Pour ce faire,
multiplier les deux chiffres par 1,10. Ainsi, pour l'exemple
\texttt{!wAb!} avec des chiffres de latitude ajoutés de 32, on multiplie
par 1,10 pour obtenir des chiffres réels ajoutés de 35,2. La position en
haute précision devient alors DDMM.mm352N.
\end{aprssideblock}

\begin{aprssideblock}{APRS Data\\Units}
\phantomsection\addcontentsline{toc}{section}{APRS Data Units}\label{unituxe9s-de-donnuxe9es-aprs}
Pour des raisons historiques, il existe un certain manque de cohérence entre
les unités de données dans les paquets APRS --- certaines vitesses sont
exprimées en knots (nœuds), d'autres en miles per hour (miles par heure) ;
certaines altitudes sont exprimées en pieds, d'autres en mètres, et ainsi de
suite. Il est souligné que cette spécification décrit les unités de données
telles qu'elles sont transmises sur l'air (\emph{on-air}). Il incombe aux
applications APRS de convertir les unités on-air en unités plus appropriées
si nécessaire.

Le datum terrestre GPS par défaut est le World Geodetic System (WGS) 1984,
mais des options continentales telles que OSG pour le Royaume-Uni sont
acceptables.
\end{aprssideblock}
